\chapter{Red convolucional y optimizadores}
\label{cap:cnn}

Los modelos convexos del Capítulo~\ref{cap:convexos} trabajan con 90 números por ventana: la potencia de cada canal
en cinco bandas fijas. Una red convolucional (CNN) puede trabajar directamente con el espectrograma completo y aprender
qué patrones de frecuencia y tiempo indican una crisis. El precio es que su entrenamiento deja de ser un problema convexo
(Sección~\ref{sec:no_convexo}): el resultado depende del punto inicial, del algoritmo de optimización y de la tasa de
aprendizaje. Este capítulo describe la red, deduce cómo se calcula su gradiente, presenta los seis optimizadores que se
comparan y detalla el diseño del experimento. Los resultados están en el Capítulo~\ref{cap:resultados}. El código es
\texttt{scripts/entrenar\_cnn.py} y el análisis, \texttt{etapa4\_cnn\_optimizadores.ipynb}.

\section{Entrada}
\label{sec:cnn_entrada}

Cada ventana de 4~s se representa con el tensor que se guardó en el preprocesamiento (Sección~\ref{sec:preproc}):
\begin{equation}
  \mathbf{X} \in \R^{18 \times 40 \times 7}, \qquad
  X_{c,f,t} = \log_{10}\big|\mathrm{STFT}\{x_c\}(f, t)\big|^2,
\end{equation}
con 18 canales bipolares $c$, 40 frecuencias $f$ de 1 a 40~Hz (resolución de 1~Hz, ventana de Hann de 1~s) y 7 cuadros de
tiempo $t$ (salto de 0,5~s). Antes de entrar a la red se aplican dos transformaciones:
\begin{enumerate}
  \item \textbf{Calibración por paciente} (ecuación~\eqref{eq:calibracion}): a cada canal y frecuencia se le resta la mediana de
        ese paciente sobre todas sus ventanas y cuadros. Elimina el término $2\log g_{p,c}$ del modelo~\eqref{eq:modelo_eeg}
        y, además, la forma espectral basal propia de cada paciente.
  \item \textbf{Estandarización}: a cada canal y frecuencia se le resta la media y se lo divide por el desvío, calculados
        solo con el conjunto de entrenamiento.
\end{enumerate}
Los canales del EEG se tratan como los \emph{canales de entrada} de la convolución (igual que los tres colores de una
imagen), y el plano frecuencia--tiempo, como la ``imagen''. Así cada filtro de la primera capa combina los 18 canales
con pesos propios, y se desplaza a lo largo de la frecuencia y del tiempo.

\section{Arquitectura}
\label{sec:arquitectura}

\subsection{Capas}

\paragraph{Convolución.} Una capa convolucional con $C_\text{in}$ canales de entrada, $C_\text{out}$ de salida y núcleo de
$k_f \times k_t$ calcula
\begin{equation}
  Z_{o,f,t} = b_o + \sum_{c=1}^{C_\text{in}}\ \sum_{i=1}^{k_f}\ \sum_{j=1}^{k_t} W_{o,c,i,j}\; X_{c,\ f + i - r_f,\ t + j - r_t},
  \qquad r_f = \tfrac{k_f + 1}{2},\ r_t = \tfrac{k_t + 1}{2},
  \label{eq:conv}
\end{equation}
con relleno de ceros en los bordes para que la salida tenga el mismo tamaño que la entrada. Estrictamente es una
\emph{correlación cruzada} (la convolución del Capítulo~\ref{cap:fourier} invierte el núcleo), pero como $W$ se aprende la
diferencia es irrelevante. Tiene $C_\text{out}(C_\text{in}k_fk_t + 1)$ parámetros, independientemente del tamaño de la entrada:
los mismos pesos se usan en todas las posiciones.

Esa repetición de pesos da la propiedad que motiva la arquitectura: la capa es \textbf{equivariante a traslaciones}. Si la
entrada se desplaza $\Delta$ posiciones a lo largo de la frecuencia, la salida se desplaza igual. Una descarga rítmica con
fundamental $f_0$ (ecuación~\eqref{eq:descarga}) aparece como una línea horizontal en el espectrograma; si en otro paciente
$f_0$ es distinta, la línea está más arriba o más abajo, y el mismo filtro la detecta. Es la variabilidad entre pacientes que
mostró la Tabla~\ref{tab:dif_log}. (Como el eje de frecuencias es lineal, los armónicos $2f_0, 3f_0$ no se desplazan lo mismo
que la fundamental, así que la equivariancia es solo aproximada para el patrón completo.)

\paragraph{Normalización por lotes \parencite{ioffe2015}.} Para cada canal $o$, sobre las $B$ ventanas del lote y todas
las posiciones $(f, t)$, se calcula la media $\mu_o$ y la varianza $\sigma_o^2$, y
\begin{equation}
  \hat Z_{o,f,t} = \frac{Z_{o,f,t} - \mu_o}{\sqrt{\sigma_o^2 + \epsilon}}, \qquad
  \text{BN}(Z)_{o,f,t} = \gamma_o\,\hat Z_{o,f,t} + \beta_o,
  \label{eq:bn}
\end{equation}
con $\gamma_o, \beta_o$ aprendidos y $\epsilon = 10^{-5}$. Durante la evaluación no hay lote: se usan promedios móviles de
$\mu_o$ y $\sigma_o^2$ acumulados en el entrenamiento (factor 0,1 por paso). Su efecto sobre la optimización se analiza en la
Sección~\ref{sec:bn_escala}.

\paragraph{ReLU} $\max(0, z)$, elemento a elemento. \textbf{Max pooling} $(2,1)$: toma el máximo de cada par de frecuencias
vecinas, reduciendo a la mitad el eje de frecuencia y dejando el tiempo igual. \textbf{Promedio global}: promedia cada
canal sobre todas las posiciones, dejando un vector de 64 números. \textbf{Dropout} \parencite{goodfellow2016}: durante el
entrenamiento pone en cero cada componente con probabilidad $p = 0{,}3$ y multiplica las restantes por $1/(1-p)$, de modo que
la esperanza no cambia; en evaluación no hace nada. Finalmente una capa lineal da un único número, el \emph{logit} $z$.

\subsection{Tabla de capas}

\begin{table}[ht]
\centering
\caption{Arquitectura de la CNN. Entrada: $18 \times 40 \times 7$ (canales $\times$ frecuencias $\times$ cuadros de tiempo).
Los parámetros incluyen pesos, sesgos y los $\gamma, \beta$ de la normalización por lotes.}
\label{tab:cnn_capas}
\begin{tabular}{clcr}
\toprule
\# & Bloque & Salida & Parámetros \\
\midrule
1 & Conv $5\times3$, 32 filtros + BN + ReLU + MaxPool $(2,1)$ & $32 \times 20 \times 7$ & 8\,736 \\
2 & Conv $3\times3$, 64 filtros + BN + ReLU + MaxPool $(2,1)$ & $64 \times 10 \times 7$ & 18\,624 \\
3 & Conv $3\times3$, 64 filtros + BN + ReLU                   & $64 \times 10 \times 7$ & 37\,056 \\
4 & Promedio global + Dropout 0,3 + Lineal                    & 1 (logit)               & 65 \\
\midrule
  & Total &  & 64\,481 \\
\bottomrule
\end{tabular}
\end{table}

Por ejemplo, el bloque 1 tiene $32\cdot(18\cdot5\cdot3 + 1) = 8\,672$ parámetros de convolución y $2\cdot32 = 64$ de la
normalización. El primer núcleo es más alto en frecuencia ($5$) que en tiempo ($3$) para captar picos espectrales de
algunos hertz de ancho.

\paragraph{Campo receptivo.} Cada salida de la última capa convolucional depende de una región de la entrada que se calcula
capa por capa: un núcleo de tamaño $k$ agrega $(k-1)\cdot s$ posiciones, donde $s$ es el producto de los pasos de los
\emph{poolings} anteriores. En frecuencia: $5 \to 6 \to 10 \to 12 \to 20$, es decir 20~Hz. En tiempo: $3 \to 5 \to 7$
cuadros, la ventana entera. Cada neurona de la última capa ve la mitad del rango de frecuencias y los 4~s completos, y el
promedio global combina todas.

\paragraph{Tamaño.} 64\,481 parámetros es poco para una CNN, y es deliberado: hay solo 403 ventanas ictales de
entrenamiento (13 crisis de 3 pacientes), y una red más grande las memorizaría. Para comparar, la regresión logística de
CHB-MIT tiene 91 parámetros.

\subsection{Salida y pérdida}

La salida $z = f_{\vth}(\mathbf{X})$ se interpreta como el logit de la probabilidad de crisis, $\hat p = \sigma(z)$, y la red
se entrena minimizando la entropía cruzada binaria promedio sobre el lote $\mathcal{B}$:
\begin{equation}
  J_{\mathcal{B}}(\vth) = \frac{1}{|\mathcal{B}|}\sum_{i \in \mathcal{B}} \ell(z_i, y_i), \qquad
  \ell(z, y) = -y\log\sigma(z) - (1 - y)\log\big(1 - \sigma(z)\big).
  \label{eq:bce}
\end{equation}
Es la pérdida~\eqref{eq:logistica_binaria} de la regresión logística, con la diferencia de que $z$ ya no es lineal en $\vth$.
Por eso es convexa en $z$ pero no en $\vth$. Numéricamente se calcula como
$\ell = \max(z, 0) - zy + \log(1 + e^{-|z|})$, que no desborda para $|z|$ grande (función
\texttt{binary\_cross\_entropy\_with\_logits}). Como la decisión final usa un umbral de Neyman-Pearson sobre $z$
(Sección~\ref{sec:umbral_np}), la calibración de $\hat p$ como probabilidad no importa, solo el orden de los puntajes.

\section{Retropropagación}
\label{sec:retropropagacion}

Todos los optimizadores de este capítulo necesitan el gradiente $\nabla_{\vth}J_{\mathcal{B}}$. La red es una composición
de funciones, $z = h_L \circ h_{L-1} \circ \cdots \circ h_1(\mathbf{X})$, cada una con sus parámetros $\vth_\ell$. La
\textbf{retropropagación} es la regla de la cadena aplicada de la salida hacia la entrada, reutilizando los resultados
intermedios.

\paragraph{Último paso.} Con $\sigma'(z) = \sigma(z)(1 - \sigma(z))$,
\begin{equation}
  \frac{\partial\ell}{\partial z} = -y\,\frac{\sigma(z)(1 - \sigma(z))}{\sigma(z)} + (1 - y)\,\frac{\sigma(z)(1 - \sigma(z))}{1 - \sigma(z)}
  = \sigma(z) - y = \hat p - y.
  \label{eq:grad_bce}
\end{equation}
Es el caso binario de~\eqref{eq:grad_ce}: \emph{predicción menos verdad}. El gradiente es grande cuando la red se equivoca
con confianza y casi cero cuando acierta con confianza: las ventanas fáciles casi no aportan al aprendizaje.

\paragraph{Recursión.} Sea $\mathbf{a}_\ell = h_\ell(\mathbf{a}_{\ell-1}; \vth_\ell)$ la salida de la capa $\ell$ y
$\boldsymbol\delta_\ell = \partial\ell/\partial\mathbf{a}_\ell$ (el ``error'' que llega a esa capa). Partiendo de
$\delta_L = \hat p - y$,
\begin{equation}
  \frac{\partial\ell}{\partial\vth_\ell} = \Big(\frac{\partial\mathbf{a}_\ell}{\partial\vth_\ell}\Big)^{\!\top}\boldsymbol\delta_\ell,
  \qquad
  \boldsymbol\delta_{\ell-1} = \Big(\frac{\partial\mathbf{a}_\ell}{\partial\mathbf{a}_{\ell-1}}\Big)^{\!\top}\boldsymbol\delta_\ell.
  \label{eq:backprop}
\end{equation}
Nunca se forman las matrices jacobianas: cada capa sabe calcular directamente el producto de su jacobiana transpuesta por un
vector. Para las capas usadas:
\begin{itemize}
  \item \textbf{Lineal} $z = \vw^\top\mathbf{a} + b$: $\partial\ell/\partial\vw = \delta\,\mathbf{a}$, $\partial\ell/\partial b = \delta$,
        $\boldsymbol\delta_\text{entrada} = \delta\,\vw$.
  \item \textbf{Convolución}~\eqref{eq:conv}: derivando,
        \[
          \frac{\partial\ell}{\partial W_{o,c,i,j}} = \sum_{f,t}\delta_{o,f,t}\,X_{c,\,f+i-r_f,\,t+j-r_t},
          \qquad
          \frac{\partial\ell}{\partial b_o} = \sum_{f,t}\delta_{o,f,t}.
        \]
        El gradiente de cada peso suma la contribución de todas las posiciones donde se usó; el error hacia atrás es
        otra convolución, con el núcleo invertido.
  \item \textbf{ReLU}: deja pasar $\delta$ donde la entrada fue positiva y lo anula donde no.
  \item \textbf{Max pooling}: el error va entero a la posición que tuvo el máximo; las demás reciben cero.
  \item \textbf{Promedio global} sobre $P$ posiciones: cada posición recibe $\delta/P$.
  \item \textbf{Normalización por lotes}: como $\mu_o$ y $\sigma_o$ dependen de todo el lote, el gradiente de una ventana
        depende de las otras. Llamando $\mathbf{g} = \partial\ell/\partial\hat Z$ (sobre las $m$ entradas del canal),
        $\partial\ell/\partial Z = \frac{1}{\sqrt{\sigma^2 + \epsilon}}\big(\mathbf{g} - \bar g - \hat Z\,\overline{g\hat Z}\big)$,
        donde la barra es el promedio sobre las $m$ entradas: es $\mathbf{g}$ proyectado para quitarle las componentes en la
        dirección de la media y de la escala.
\end{itemize}
Calcular el gradiente cuesta del orden de dos veces lo que cuesta evaluar la red, independientemente del número de
parámetros. En la práctica lo calcula la diferenciación automática de PyTorch (\texttt{loss.backward()}), que implementa
exactamente~\eqref{eq:backprop}.

\subsection{Invariancia de escala por la normalización}
\label{sec:bn_escala}

La normalización por lotes cambia la geometría del problema de optimización de una forma que importa para comparar
optimizadores.

\begin{proposicion}
\label{prop:bn_escala}
Sea $\vw$ el vector de pesos de una convolución seguida de normalización por lotes, y $J(\vw)$ la pérdida vista como
función de $\vw$ (con el resto fijo). Si se desprecia $\epsilon$, para todo $c > 0$:
\begin{enumerate}
  \item $J(c\vw) = J(\vw)$;
  \item $\nabla J(c\vw) = \frac{1}{c}\nabla J(\vw)$;
  \item $\vw^\top\nabla J(\vw) = 0$.
\end{enumerate}
\end{proposicion}
\begin{proof}
(1) Multiplicar $\vw$ por $c$ multiplica la salida de la convolución (sin el sesgo, que la normalización cancela) por $c$, y
con ella la media y el desvío del lote; el cociente~\eqref{eq:bn} no cambia. (2) Derivando $J(c\vw) = J(\vw)$ respecto de
$\vw$: $c\,\nabla J(c\vw) = \nabla J(\vw)$. (3) Derivando $J(c\vw) = J(\vw)$ respecto de $c$ en $c = 1$:
$\vw^\top\nabla J(\vw) = 0$.
\end{proof}

Las consecuencias son dos. Primero, la pérdida \emph{no tiene mínimo aislado} en esos pesos: toda la semirrecta
$\{c\vw^* : c > 0\}$ tiene el mismo valor. Es otra forma de no convexidad estricta, además de las simetrías de la
Sección~\ref{sec:no_convexo}. Segundo, como el gradiente es ortogonal a $\vw$, un paso de SGD aumenta la norma:
$\lVert\vw - \eta\nabla J\rVert^2 = \lVert\vw\rVert^2 + \eta^2\lVert\nabla J\rVert^2$; y por (2) el paso \emph{angular}, que es lo
único que cambia la función, es del orden de $\eta\lVert\nabla J(\vw)\rVert/\lVert\vw\rVert \propto \eta/\lVert\vw\rVert^2$. La tasa de
aprendizaje efectiva de esas capas depende de la norma de sus pesos, que a su vez depende del optimizador y del decaimiento
de pesos. Esto hace que comparar valores de $\eta$ entre optimizadores sea todavía menos directo, y es una razón más para
buscar la tasa por separado para cada uno.

\section{Optimizadores de primer orden}
\label{sec:optimizadores}

Notación: $\mathbf{g}_t = \nabla J_{\mathcal{B}_t}(\vth_t)$ es el gradiente sobre el lote $\mathcal{B}_t$ de 64 ventanas, $\eta$ la
tasa de aprendizaje y $\lambda$ el decaimiento de pesos. Las operaciones entre vectores ($\mathbf{g}^2$, $\sqrt{\cdot}$, la
división) son elemento a elemento. Las fórmulas son las que implementa PyTorch, que en algunos casos difieren en detalles de
los artículos originales; se señalan las diferencias.

\subsection{SGD}

\begin{equation}
  \vth_{t+1} = \vth_t - \eta\,(\mathbf{g}_t + \lambda\vth_t).
  \label{eq:sgd}
\end{equation}
Es el descenso por gradiente estocástico de la Sección~\ref{sec:gradiente}. Sumar $\lambda\vth$ al gradiente equivale a
minimizar $J + \frac{\lambda}{2}\lVert\vth\rVert^2$ (regularización L2), y también a multiplicar los pesos por $(1 - \eta\lambda)$
antes de cada paso (\emph{decaimiento de pesos}); en SGD las dos cosas son lo mismo. Se usa $\lambda = 10^{-4}$.

Por el análisis de la Sección~\ref{sec:gradiente}, en una dirección de curvatura $L$ el paso debe cumplir $\eta < 2/L$ para no
diverger, y la velocidad de convergencia está limitada por la dirección más plana: con el mejor paso fijo,
$\eta = 2/(L + \mu)$, el error se reduce en un factor $(\kappa - 1)/(\kappa + 1)$ por iteración en una cuadrática con número de
condición $\kappa$.

\subsection{Momentum \parencite{polyak1964}}

El método de la ``bola pesada'' agrega una velocidad que acumula los gradientes pasados:
\begin{equation}
  \mathbf{v}_{t+1} = \mu\,\mathbf{v}_t + \mathbf{g}_t, \qquad \vth_{t+1} = \vth_t - \eta\,\mathbf{v}_{t+1},
  \qquad \mu = 0{,}9.
  \label{eq:momentum}
\end{equation}
Desarrollando, $\mathbf{v}_{t+1} = \sum_{k \ge 0}\mu^k\mathbf{g}_{t-k}$: un promedio de los gradientes con peso exponencial (memoria
de unos $1/(1-\mu) = 10$ pasos). Eliminando $\mathbf{v}$ se obtiene la forma clásica
\begin{equation}
  \vth_{t+1} = \vth_t - \eta\,\mathbf{g}_t + \mu\,(\vth_t - \vth_{t-1}).
  \label{eq:heavy_ball}
\end{equation}
Dos efectos explican por qué ayuda. En las direcciones donde el gradiente cambia de signo de un paso a otro (las paredes de
un valle angosto) los términos se cancelan y la oscilación se amortigua. En las direcciones donde el gradiente es
consistente (el fondo del valle) se suman, y el paso crece.

\paragraph{Paso efectivo.} Si el gradiente es aproximadamente constante, $\mathbf{v}$ converge a $\mathbf{g}/(1 - \mu)$ y el paso es
\begin{equation}
  \eta_\text{ef} = \frac{\eta}{1 - \mu} = 10\,\eta.
  \label{eq:paso_efectivo}
\end{equation}
Por eso las tasas de Momentum y Nesterov se buscan en una grilla diez veces menor que la de SGD
(Tabla~\ref{tab:tasas}), y las tasas se comparan entre estos tres métodos por su paso efectivo.

\paragraph{Análisis en la cuadrática.} Sea $J(\vth) = \frac12\vth^\top A\vth$ con $A$ simétrica y autovalores en $[\mu_A, L]$,
$\kappa = L/\mu_A$, y gradiente exacto. En la base de autovectores las coordenadas se desacoplan, y en una dirección de
autovalor $a$ la recursión~\eqref{eq:heavy_ball} es lineal de segundo orden:
\begin{equation}
  x_{t+1} = (1 + \mu - \eta a)\,x_t - \mu\,x_{t-1}.
\end{equation}

\begin{proposicion}
\label{prop:momentum}
\begin{enumerate}
  \item La recursión converge a cero si y solo si $0 < \eta a < 2(1 + \mu)$, con $0 \le \mu < 1$.
  \item Si $(1 - \sqrt\mu)^2 \le \eta a \le (1 + \sqrt\mu)^2$, el factor de convergencia es exactamente $\sqrt\mu$.
  \item Con $\eta = \dfrac{4}{(\sqrt L + \sqrt{\mu_A})^2}$ y $\sqrt\mu = \dfrac{\sqrt\kappa - 1}{\sqrt\kappa + 1}$, todas las direcciones
        convergen con factor
        \begin{equation}
          \frac{\sqrt\kappa - 1}{\sqrt\kappa + 1}
          \qquad\text{frente a}\qquad
          \frac{\kappa - 1}{\kappa + 1}
          \ \text{del descenso por gradiente con el mejor paso.}
          \label{eq:tasa_momentum}
        \end{equation}
\end{enumerate}
\end{proposicion}
\begin{proof}
Las soluciones son $x_t = c_1r_1^t + c_2r_2^t$, con $r_{1,2}$ raíces de $p(r) = r^2 - (1 + \mu - \eta a)\,r + \mu$. Convergen si
y solo si $|r_{1,2}| < 1$.

(1) Por el criterio de Jury, un polinomio $r^2 + c_1r + c_0$ tiene las dos raíces en el disco unidad si y solo si $|c_0| < 1$,
$p(1) = 1 + c_1 + c_0 > 0$ y $p(-1) = 1 - c_1 + c_0 > 0$. Aquí $c_0 = \mu$, $p(1) = \eta a$ y $p(-1) = 2(1 + \mu) - \eta a$.

(2) Las raíces son complejas conjugadas cuando el discriminante es negativo, $(1 + \mu - \eta a)^2 \le 4\mu$, es decir
$|1 + \mu - \eta a| \le 2\sqrt\mu$, que equivale a $(1 - \sqrt\mu)^2 \le \eta a \le (1 + \sqrt\mu)^2$. En ese caso
$|r_1|^2 = |r_2|^2 = r_1r_2 = \mu$, el término independiente: el módulo es $\sqrt\mu$ sin importar $a$.

(3) Se eligen $\eta$ y $\mu$ para que el intervalo de (2) sea exactamente $[\eta\mu_A, \eta L]$:
$1 - \sqrt\mu = \sqrt{\eta\mu_A}$ y $1 + \sqrt\mu = \sqrt{\eta L}$. Sumando, $2 = \sqrt\eta(\sqrt L + \sqrt{\mu_A})$; restando,
$2\sqrt\mu = \sqrt\eta(\sqrt L - \sqrt{\mu_A})$. Dividiendo, $\sqrt\mu = (\sqrt L - \sqrt{\mu_A})/(\sqrt L + \sqrt{\mu_A})$, y dividiendo
numerador y denominador por $\sqrt{\mu_A}$ se obtiene~\eqref{eq:tasa_momentum}.

Para el descenso por gradiente ($\mu = 0$) el factor en la dirección $a$ es $|1 - \eta a|$; el peor caso sobre
$a \in [\mu_A, L]$ se minimiza igualando los extremos, $1 - \eta\mu_A = \eta L - 1$, lo que da $\eta = 2/(L + \mu_A)$ y factor
$(L - \mu_A)/(L + \mu_A) = (\kappa - 1)/(\kappa + 1)$.
\end{proof}

La ganancia está en reemplazar $\kappa$ por $\sqrt\kappa$. Para reducir el error en un factor $\epsilon$ hacen falta del orden de
$\frac{\kappa}{2}\log\frac1\epsilon$ iteraciones sin momento y $\frac{\sqrt\kappa}{2}\log\frac1\epsilon$ con momento. Con el
$\kappa \approx 528$ de la regresión logística de Bonn (Sección~\ref{sec:gradiente}): los factores son $0{,}9962$ y $0{,}9166$, y
para reducir la distancia al óptimo en 12 órdenes de magnitud hacen falta unas 7\,300 iteraciones sin momento y unas 320 con
momento.

El punto (1) también es informativo: con $\mu = 0{,}9$ la región estable llega hasta $\eta a < 3{,}8$, mientras que sin momento
llega hasta $\eta a < 2$. Medido en paso efectivo~\eqref{eq:paso_efectivo}, la condición con momento es
$\eta_\text{ef}\,a < 38$: a igual paso efectivo, el momento tolera curvaturas mucho mayores. Este análisis vale para una
cuadrática con gradiente exacto; con la pérdida no convexa de la CNN y gradientes estocásticos es solo una guía. En
particular, el $\mu$ óptimo depende de $\kappa$, que no se conoce, y en la práctica se fija $\mu = 0{,}9$.

\subsection{Nesterov \parencite{nesterov1983,sutskever2013}}

El gradiente acelerado de Nesterov evalúa el gradiente en el punto al que el momento va a llevar los parámetros, no en el
punto actual:
\begin{equation}
  \mathbf{v}_{t+1} = \mu\,\mathbf{v}_t + \nabla J(\vth_t - \eta\mu\,\mathbf{v}_t), \qquad \vth_{t+1} = \vth_t - \eta\,\mathbf{v}_{t+1}.
  \label{eq:nesterov}
\end{equation}
Si el momento está por pasarse del mínimo, el gradiente en el punto adelantado ya apunta hacia atrás y frena antes. Para
funciones convexas $L$-suaves, la versión original alcanza un error $O(1/k^2)$ en lugar del $O(1/k)$ del
Teorema~\ref{teo:convergencia}, que es la mejor tasa posible para métodos que solo usan gradientes; para funciones
$\mu$-fuertemente convexas el factor por iteración es $1 - 1/\sqrt\kappa$.

PyTorch usa una forma equivalente que evita calcular el gradiente en otro punto, cambiando de variable al punto adelantado
\parencite{sutskever2013}:
\begin{equation}
  \mathbf{v}_{t+1} = \mu\,\mathbf{v}_t + \mathbf{g}_t, \qquad \vth_{t+1} = \vth_t - \eta\,(\mathbf{g}_t + \mu\,\mathbf{v}_{t+1}).
\end{equation}
Comparada con~\eqref{eq:momentum}, el paso usa el gradiente actual con peso $1 + \mu$ y la velocidad vieja con peso $\mu^2$:
reacciona más rápido a un cambio de dirección. El paso efectivo con gradiente constante es el mismo, $\eta/(1 - \mu)$.

\subsection{RMSprop \parencite{tieleman2012}}

Los métodos anteriores usan la misma tasa para todos los parámetros. Pero en una red los gradientes de distintas capas
pueden diferir en órdenes de magnitud (Sección~\ref{sec:bn_escala}). RMSprop divide cada coordenada por una estimación de la
magnitud típica de su gradiente:
\begin{equation}
  \mathbf{s}_{t+1} = \rho\,\mathbf{s}_t + (1 - \rho)\,\mathbf{g}_t^2, \qquad
  \vth_{t+1} = \vth_t - \eta\,\frac{\mathbf{g}_t}{\sqrt{\mathbf{s}_{t+1}} + \epsilon},
  \qquad \rho = 0{,}99,\ \epsilon = 10^{-8}.
  \label{eq:rmsprop}
\end{equation}
$\mathbf{s}$ es un promedio exponencial de los cuadrados, con memoria de unos 100 pasos, y $\sqrt{\mathbf{s}}$ es la raíz cuadrática
media (RMS) reciente de cada coordenada. Si el gradiente de una coordenada tiene magnitud estable $|g|$, el paso en esa
coordenada es $\approx\eta\,\mathrm{signo}(g)$: \textbf{la tasa queda medida en unidades del parámetro}, no del gradiente. Por eso
los métodos adaptativos usan tasas mucho más chicas que SGD ($10^{-3}$ en lugar de $10^{-1}$) y por eso sus tasas no se
pueden comparar con las de SGD.

Formalmente es un \emph{precondicionamiento diagonal}: un paso de descenso por gradiente con la métrica
$D = \operatorname{diag}(\sqrt{\mathbf{s}} + \epsilon)$, $\vth_{t+1} = \vth_t - \eta D^{-1}\mathbf{g}_t$. Si la Hessiana fuera diagonal y
$\sqrt{\mathbf{s}}$ fuera proporcional a sus entradas, $D^{-1}$ la llevaría a la identidad y el número de condición a 1. En general
la Hessiana no es diagonal y la corrección es parcial.

Un detalle de la implementación: PyTorch inicia $\mathbf{s}_0 = \mathbf{0}$ y no corrige el sesgo. En el primer paso
$\mathbf{s}_1 = 0{,}01\,\mathbf{g}_0^2$ y el paso es $\eta\,\mathbf{g}_0/(0{,}1|\mathbf{g}_0|) = 10\,\eta$ por coordenada. Los primeros pasos de RMSprop son
entonces hasta diez veces más grandes que los siguientes, hasta que $\mathbf{s}$ se llena (unos 100 pasos, cerca de
una época y media).

\subsection{Adam \parencite{kingma2015}}

Adam combina el momento (un promedio exponencial del gradiente) con la normalización de RMSprop, y corrige el sesgo de
inicializar los promedios en cero:
\begin{align}
  \mathbf{m}_{t+1} &= \beta_1\mathbf{m}_t + (1 - \beta_1)\,\mathbf{g}_t, &
  \mathbf{s}_{t+1} &= \beta_2\mathbf{s}_t + (1 - \beta_2)\,\mathbf{g}_t^2, \nonumber\\
  \hat{\mathbf{m}} &= \frac{\mathbf{m}_{t+1}}{1 - \beta_1^{t+1}}, &
  \hat{\mathbf{s}} &= \frac{\mathbf{s}_{t+1}}{1 - \beta_2^{t+1}}, \label{eq:adam}\\
  \vth_{t+1} &= \vth_t - \eta\,\frac{\hat{\mathbf{m}}}{\sqrt{\hat{\mathbf{s}}} + \epsilon}, &
  &\beta_1 = 0{,}9,\ \beta_2 = 0{,}999,\ \epsilon = 10^{-8}. \nonumber
\end{align}

\paragraph{Corrección del sesgo.} Desarrollando, $\mathbf{m}_t = (1 - \beta_1)\sum_{k=0}^{t-1}\beta_1^{k}\mathbf{g}_{t-1-k}$. Si los
gradientes tuvieran esperanza constante $\bar{\mathbf{g}}$,
\[
  \E[\mathbf{m}_t] = (1 - \beta_1)\sum_{k=0}^{t-1}\beta_1^k\,\bar{\mathbf{g}} = (1 - \beta_1^t)\,\bar{\mathbf{g}},
\]
y dividir por $1 - \beta_1^t$ da un estimador insesgado. Lo mismo para $\mathbf{s}$ con $\beta_2$. Sin corregir, con
$\beta_2 = 0{,}999$ el denominador sería casi cero en los primeros cientos de pasos y los pasos, enormes (el problema que se
señaló en RMSprop, pero cien veces peor).

\paragraph{Tamaño del paso.} Si en una coordenada el gradiente es consistente, $\hat m \approx \sqrt{\hat s}$ y el paso es
$\approx\eta$. Si es ruido alrededor de cero, $\hat m$ es chico frente a $\sqrt{\hat s}$ y el paso se achica. El cociente
$\hat{\mathbf{m}}/\sqrt{\hat{\mathbf{s}}}$ funciona como una relación señal/ruido por coordenada. El paso también es invariante a
reescalar la pérdida por una constante, porque el factor se cancela entre numerador y denominador.

\paragraph{Decaimiento de pesos acoplado.} PyTorch implementa el \texttt{weight\_decay} de Adam sumándolo al gradiente,
$\mathbf{g}_t \leftarrow \mathbf{g}_t + \lambda\vth_t$, \emph{antes} de calcular $\mathbf{m}$ y $\mathbf{s}$. Se usa $\lambda = 10^{-4}$.

\subsection{AdamW \parencite{loshchilov2019}}

En SGD, sumar $\lambda\vth$ al gradiente y achicar los pesos son lo mismo. En Adam no: el término $\lambda\vth$ entra en
$\hat{\mathbf{m}}$ y después se divide por $\sqrt{\hat{\mathbf{s}}}$. Si se ignora el momento, la parte del paso que viene de la
regularización en la coordenada $j$ es
\begin{equation}
  \eta\,\frac{\lambda\,\theta_j}{\sqrt{\hat s_j} + \epsilon}.
\end{equation}
Los parámetros con gradientes grandes ($\hat s_j$ grande) casi no se regularizan, y los que tienen gradientes chicos se
regularizan mucho. El resultado ya no corresponde a minimizar $J + \frac{\lambda}{2}\lVert\vth\rVert^2$, ni a ninguna penalización
fija, porque $\hat{\mathbf{s}}$ cambia durante el entrenamiento. AdamW \emph{desacopla} el decaimiento y lo aplica directamente
sobre los pesos:
\begin{equation}
  \vth_{t+1} = (1 - \eta\lambda)\,\vth_t - \eta\,\frac{\hat{\mathbf{m}}}{\sqrt{\hat{\mathbf{s}}} + \epsilon},
  \label{eq:adamw}
\end{equation}
con $\hat{\mathbf{m}}$ y $\hat{\mathbf{s}}$ calculados con el gradiente de $J$ solo. Ahora todos los pesos se achican en la misma
proporción. Como el factor $\eta\lambda$ es mucho menor que en SGD (porque $\eta$ es chico), se usa un $\lambda$ mayor:
$\lambda = 10^{-2}$, el valor por defecto de PyTorch.

En este experimento el efecto es pequeño. Con $\eta = 3\times10^{-4}$, cada paso multiplica los pesos por
$1 - 3\times10^{-6}$; en unos 2100 pasos (Sección~\ref{sec:diseno}) eso es un factor $0{,}994$. La diferencia entre Adam y AdamW
se reduce a esa leve contracción y a la forma en que se regularizan los parámetros con gradientes chicos.

\subsection{Resumen}

\begin{table}[ht]
\centering
\caption{Los seis optimizadores comparados. ``Paso'' indica qué controla $\eta$: un múltiplo del gradiente o una magnitud en
unidades del parámetro.}
\label{tab:optimizadores}
\small
\begin{tabular}{lllll}
\toprule
Optimizador & Momento & Adaptativo & Decaimiento & Paso \\
\midrule
SGD      & no                    & no & L2, $\lambda = 10^{-4}$          & $\eta\,\mathbf{g}$ \\
Momentum & $\mu = 0{,}9$         & no & L2, $\lambda = 10^{-4}$          & $\approx 10\,\eta\,\mathbf{g}$ \\
Nesterov & $\mu = 0{,}9$, adelantado & no & L2, $\lambda = 10^{-4}$      & $\approx 10\,\eta\,\mathbf{g}$ \\
RMSprop  & no                    & $\rho = 0{,}99$ & acoplado, $10^{-4}$ & $\approx\eta$ por coordenada \\
Adam     & $\beta_1 = 0{,}9$     & $\beta_2 = 0{,}999$ & acoplado, $10^{-4}$ & $\approx\eta$ por coordenada \\
AdamW    & $\beta_1 = 0{,}9$     & $\beta_2 = 0{,}999$ & desacoplado, $10^{-2}$ & $\approx\eta$ por coordenada \\
\bottomrule
\end{tabular}
\end{table}

El decaimiento de pesos se aplica a todos los parámetros, incluidos los sesgos y los $\gamma, \beta$ de la normalización por
lotes.

\subsection{Qué dice la teoría sobre la generalización}

Las garantías de esta sección son sobre la \emph{pérdida de entrenamiento} y en problemas convexos. En una red, lo que
interesa es el desempeño sobre pacientes que no se usaron para entrenar, y para eso no hay teoremas comparables. Hay dos
observaciones empíricas que sirven para interpretar los resultados:
\begin{itemize}
  \item \textcite{wilson2017} encontraron, en varios problemas de clasificación, que los métodos adaptativos (RMSprop, Adam)
        llegan a pérdidas de entrenamiento iguales o menores que SGD pero generalizan peor, y dieron un ejemplo construido
        donde la solución a la que converge Adam es peor que la de SGD.
  \item Con pasos grandes, el ruido del gradiente estocástico impide que el optimizador se asiente en mínimos angostos, donde
        la pérdida sube rápido al moverse un poco, y lo empuja hacia regiones más planas, que suelen generalizar mejor
        \parencite{goodfellow2016}. Según esta idea, lo que importa no es tanto el algoritmo como el tamaño del paso, y la
        relación entre la tasa de aprendizaje y el tamaño del lote.
\end{itemize}
El Capítulo~\ref{cap:resultados} pone a prueba estas dos ideas en este problema.

\section{Diseño experimental}
\label{sec:diseno}

\subsection{Datos por época}

Con la partición fija (Tabla~\ref{tab:particion}), el entrenamiento tiene 403 ventanas ictales y unas 84\,400 sin crisis
(excluidas las de borde). Usar todas en cada época sería lento y la pérdida quedaría dominada por las ventanas normales. En
cada época se toman:
\begin{itemize}
  \item las 403 ventanas ictales;
  \item $10 \times 403 = 4030$ ventanas sin crisis elegidas al azar sin reposición, \emph{distintas en cada época}.
\end{itemize}
Son 4433 ventanas, mezcladas y divididas en 70 lotes de 64. En 30 épocas (unos 2100 pasos de optimización) la red ve
unas 120\,000 ventanas normales, más que las que hay: cada una aparece en promedio 1,4 veces. Una de cada 11 ventanas de
entrenamiento es ictal, frente a una de cada 210 en los datos de entrenamiento. Este submuestreo solo desplaza los logits (como el peso
de clases de la Sección~\ref{sec:neyman_pearson}), no cambia el orden de los puntajes que la red busca aprender, y el umbral
se fija después con Neyman-Pearson sobre datos con la proporción real.

\subsection{Validación y elección de la época}

Después de cada época se evalúa la red sobre un subconjunto fijo del paciente de validación (chb01): sus 107 ventanas
ictales y 10\,000 ventanas sin crisis elegidas al azar una sola vez, iguales para todas las corridas. Se guarda la red de la
época con \textbf{mayor AUC de validación} (\emph{early stopping}: las 30 épocas se corren siempre, pero se conserva la mejor).
Si la pérdida o los puntajes dejan de ser números finitos, la corrida se corta y se informa como divergente. Al final, la red
guardada se evalúa sobre las 36\,496 ventanas de chb01 y las 79\,221 de prueba.

\subsection{Tasas de aprendizaje}

Cada optimizador tiene su escala natural de tasas (Sección~\ref{sec:optimizadores}), así que no tiene sentido usar la misma
grilla para todos. Se prueban tres valores por optimizador, separados por factores de aproximadamente 3, alrededor de los valores
habituales (Tabla~\ref{tab:tasas}).

\begin{table}[ht]
\centering
\caption{Tasas de aprendizaje probadas. Para Momentum y Nesterov se indica entre paréntesis el paso efectivo
$\eta/(1 - \mu)$, comparable con la tasa de SGD.}
\label{tab:tasas}
\begin{tabular}{lccc}
\toprule
Optimizador & \multicolumn{3}{c}{Tasas $\eta$} \\
\midrule
SGD               & 0,3 & 0,1 & 0,03 \\
Momentum, Nesterov & 0,03 (0,3) & 0,01 (0,1) & 0,003 (0,03) \\
RMSprop, Adam, AdamW & $3\times10^{-3}$ & $10^{-3}$ & $3\times10^{-4}$ \\
\bottomrule
\end{tabular}
\end{table}

Con la grilla elegida así, SGD, Momentum y Nesterov recorren los mismos pasos efectivos (0,03, 0,1 y 0,3).

\subsection{Protocolo}

\begin{enumerate}
  \item \textbf{Búsqueda de tasa.} Cada optimizador con sus 3 tasas y semilla 0: 18 corridas de 30 épocas. Para cada
        optimizador se elige la tasa con mayor AUC de validación.
  \item \textbf{Semillas.} Cada optimizador con su tasa elegida y 5 semillas (0 a 4; la 0 ya está de la búsqueda): 24
        corridas más. La semilla fija la inicialización de los pesos (la de PyTorch por defecto: uniforme con escala
        $1/\sqrt{\text{fan\_in}}$), las ventanas normales de cada época, el orden de los lotes y las máscaras de dropout. Se
        informan media, desvío, mínimo y máximo del AUC de prueba, y se compara cada optimizador contra el mejor con el test
        de Mann-Whitney y corrección de Bonferroni (Sección~\ref{sec:tests}).
  \item \textbf{Modelo final.} Se elige el optimizador con mayor AUC de validación medio y, dentro de él, la semilla con mayor
        AUC de validación. Ese modelo se evalúa con las métricas clínicas: suavizado causal de 3 ventanas, umbral de
        Neyman-Pearson con $\alpha = 10^{-3}$ calculado sobre las ventanas sin crisis de chb01, sensibilidad por evento,
        falsas alarmas por hora y latencia (Capítulo~\ref{cap:metricas}).
  \item \textbf{Experimento de sensibilidad} (\emph{exploratorio}). Después de ver los resultados de prueba se agregó una
        fase en la que cada optimizador se corre con la tasa más alta de su grilla y 5 semillas. Como se diseñó mirando la
        prueba, no se usa para elegir nada; sirve para entender los resultados del paso 2.
\end{enumerate}

Ninguna decisión de los pasos 1 a 3 usa el conjunto de prueba. El paso 4 se informa separado y con esa advertencia.

\subsection{Costo}

Todo se corrió en CPU (Ryzen 7 5700U, 8 núcleos, sin GPU). Cada corrida de 30 épocas tarda entre 92 y 94~s, incluida la
evaluación por época. Las 42 corridas del diseño original más las 20 del experimento de sensibilidad suman unas
1,6 horas. Los espectrogramas de las 200\,563 ventanas se cargan en memoria en media precisión (unos 2~GB). El costo
fue el que determinó usar una partición fija y no validación cruzada por paciente para la CNN, una decisión que el
Capítulo~\ref{cap:resultados} muestra que tuvo consecuencias.
